\documentclass[10pt,a4paper]{amsart}
\usepackage[margin=19mm]{geometry}
\usepackage{amsmath,amssymb,mathtools}
\usepackage[scr=rsfs,cal=euler]{mathalfa}
\usepackage{xfrac}
\usepackage[unicode,hidelinks,bookmarks=false]{hyperref}
\newcommand{\ie}{i.e.\ }
\newcommand{\cf}{cf.\ }
\newcommand{\resp}{respectively\ }
\newcommand{\Subsection}{\subsection}
\providecommand{\nobreakdash}{\nobreak\hbox{-}}
\setlength{\emergencystretch}{2em}
\multlinegap0pt
\newcommand{\abs}[1]{\left|{#1}\right|}
\usepackage{csquotes}
\usepackage{amssymb}
\usepackage{tikz}
\newtheorem{lemma}{Lemma}
\title{A note on small cap square function and decoupling estimates for the parabola}
\author{Jongchon Kim, Liang Wang, Chun Keung Yeung}
\date{Source: arXiv:2603.07934v1 (CC BY 4.0)}
\begin{document}
\maketitle

\noindent\emph{Source and licence.} A note on small cap square function and decoupling estimates for the parabola. Version arXiv:2603.07934v1 (CC BY 4.0), distributed by arXiv under the Creative Commons Attribution 4.0 International licence, http://creativecommons.org/licenses/by/4.0/, as recorded in the arXiv licence metadata for this exact version and shown on the abstract page https://arxiv.org/abs/2603.07934v1. Full licence notice: \texttt{licenses/arxiv-2603.07934-CC-BY-4.0.txt}.\par\smallskip

\noindent\textit{Editorial scope.} Counted unit: the article's lemma br-na (source0.tex lines 249--262). The article's complete original proof of it is delivered with it (source0.tex lines 264--336, one \texttt{proof} environment of 73 lines). No mathematics is written, repaired or re-derived: the statement and the proof are the article's own lines, in the article's own notation, and no retained line is substituted or re-flowed.  Retained uncounted as the source-local context the proof needs: lines 213--222.  Nothing was omitted from any retained span, so the package carries no omission record.  No retained line contains a citation, so the wrapper carries no bibliography and no bibliography entry.  No retained line contains a figure environment, an image-inclusion command or drawing code, so no image is delivered and the package declares no asset dependency.  Editorial formatting only, in addition: an AMS \texttt{amsart} preamble that restates the article's own declarations and macros used by the retained lines, namely the environments \texttt{lemma} on separate counters, together with the standard packages those lines need.  The retained environments print in this document as Lemma 1 in order of appearance, whereas the article prints the same items with the numbering of its own full document.  The complete native source of the article as distributed in the arXiv e-print for this exact version is bundled unchanged as \texttt{sources/195979/source0.tex}.
\begin{lemma}\label{inarrow}
    Let $\{a_i\}_{i=1}^n$ be a sequence of complex numbers, and let $C\geq1$ be a positive constant. Then, at least one of the following estimates holds:
    \begin{equation}
        \bigg|\sum_{i=1}^n a_i\bigg|\leq 2C n \max_{\substack{i,j:|i-j|>1}}|a_ia_j|^{\frac{1}{2}},
    \end{equation}
    and
    \begin{equation}
        \bigg|\sum_{i=1}^n a_i\bigg|\leq \left(1+C^{-1}\right)\max_{i}\bigg|\sum_{j:|i-j|\leq1}a_j\bigg|.
    \end{equation}
\end{lemma}

\begin{lemma}\label{br-na}
The following pointwise bound holds for all $x\in\mathbb{R}^2$ and $p>0$:
\begin{equation}\label{br-na:eq1}
    |f(x)| \leq C_p |\log\delta|\bigg(
        \sum_{\theta \in \mathcal{P}\bigl(\delta^{\frac{\beta}{2}}\bigr)} |f_\theta(x)|^p
        + \sum_{i=0}^{m-1} \sum_{\tau \in \mathcal{P}(K^{-i})}
          \sum_{\substack{ \tau_1,\tau_2 \in \mathcal{P}(K^{-i-1}) \\
                          \tau_1,\tau_2 \subset \widetilde{\tau} \\
                          \operatorname{dist}(\tau_1,\tau_2) \geq  K^{-i-1} }}
          \abs{f_{\tau_1}(x)f_{\tau_2}(x)}^{\frac{p}{2}}
    \bigg)^{\frac{1}{p}},
\end{equation}
where $C_p$ is a constant depending only on $p$ and $c$ is an absolute constant.
\end{lemma}

\begin{proof}
For each $\tau\in \mathcal{P}(K^{-i})$ with $0\leq i\leq m-1$, we first observe the decomposition
\begin{equation*}
    f_{\widetilde{\tau}}(x) = \sum_{\substack{ \tau_1 \in \mathcal{P}(K^{-i-1}) \\ \tau_1\subset {\widetilde{\tau}}}} f_{\tau_1}(x).
\end{equation*}
Applying Lemma \ref{inarrow} to this sum with parameters $n=K=8$ and $C=|\log\delta|$ (note that $C=|\log\delta|\geq1$, which satisfies the hypothesis of Lemma \ref{inarrow}), we conclude that at least one of the following two estimates holds:
\begin{equation}
    |f_{\widetilde{\tau}}(x)|^p \leq \left(16|\log\delta|\right)^p
        \max_{\substack{ \tau_1,\tau_2 \in \mathcal{P}(K^{-i-1}) \\
                        \tau_1,\tau_2 \subset {\widetilde{\tau}} \\
                        \operatorname{dist}(\tau_1,\tau_2) \geq  K^{-i-1} }}
        \abs{f_{\tau_1}(x)f_{\tau_2}(x)}^{\frac{p}{2}},
\end{equation}
and
\begin{equation}
    |f_{\widetilde{\tau}}(x)|^p \leq \left(1+|\log\delta|^{-1}\right)^p
        \max_{\substack{ \tau_1 \in \mathcal{P}(K^{-i-1}) \\
                        \tau_1\subset {\widetilde{\tau}}}} |f_{\widetilde{\tau_1}}(x)|^p.
\end{equation}

Using this for the sum $f = \sum_{\tau_1 \in \mathcal{P}(K^{-1})} f_{\tau_1}$, we obtain
\begin{align}\label{br-na:eq2}
    |f(x)|^p &\leq \left(16|\log\delta|\right)^p
        \max_{\substack{ \tau_1,\tau_2 \in \mathcal{P}(K^{-1})  \\
                        \operatorname{dist}(\tau_1,\tau_2) \geq  K^{-1} }}
        \abs{f_{\tau_1}(x)f_{\tau_2}(x)}^{\frac{p}{2}}
        + \left(1+|\log\delta|^{-1}\right)^p
        \max_{\substack{ \tau_1 \in \mathcal{P}(K^{-1})}} |f_{\widetilde{\tau_1}}(x)|^p.
\end{align}
Next, for any $\tau\in\mathcal{P}(K^{-1})$, we apply the same reasoning to obtain
\begin{align}\label{br-na:eq3}
    |f_{\widetilde{\tau}}(x)|^p &\leq \left(16|\log\delta|\right)^p
        \max_{\substack{ \tau_1,\tau_2 \in \mathcal{P}(K^{-2})  \\
                        \tau_1,\tau_2\subset\widetilde{\tau}\\
                        \operatorname{dist}(\tau_1,\tau_2) \geq  K^{-2} }}
        \abs{f_{\tau_1}(x)f_{\tau_2}(x)}^{\frac{p}{2}}
        + \left(1+|\log\delta|^{-1}\right)^p
        \max_{\substack{ \tau_1 \in \mathcal{P}(K^{-2})\\
                        \tau_1\subset\widetilde{\tau }}} |f_{\widetilde{\tau_1}}(x)|^p.
\end{align}
Combining inequalities \eqref{br-na:eq2} and \eqref{br-na:eq3}, we find
\begin{align*}
|f(x)|^p &\leq \sum_{i=0}^1 \left(1+|\log\delta|^{-1}\right)^{pi} \left(16|\log\delta|\right)^p
    \sum_{\tau\in\mathcal{P}(K^{-i})}
    \max_{\substack{ \tau_1,\tau_2 \in \mathcal{P}(K^{-i-1}) \\
                    \tau_1,\tau_2 \subset {\widetilde{\tau}} \\
                    \operatorname{dist}(\tau_1,\tau_2) \geq  K^{-i-1} }}
    \abs{f_{\tau_1}(x)f_{\tau_2}(x)}^{\frac{p}{2}}\\
&\quad + \left(1+|\log\delta|^{-1}\right)^{2p}
    \max_{\substack{ \tau_1 \in \mathcal{P}(K^{-2}) }} |f_{\widetilde{\tau_1}}(x)|^p.
\end{align*}
Extending this iterative process to $m$ steps, we arrive at
\begin{align*}
|f(x)|^p
& \leq \sum_{i=0}^{m-1} \left(1+|\log\delta|^{-1}\right)^{pi} \left(16|\log\delta|\right)^p
    \sum_{\tau\in\mathcal{P}(K^{-i})}
    \max_{\substack{ \tau_1,\tau_2 \in \mathcal{P}(K^{-i-1}) \\
                    \tau_1,\tau_2 \subset {\widetilde{\tau}} \\
                    \operatorname{dist}(\tau_1,\tau_2) \geq  K^{-i-1} }}
    \abs{f_{\tau_1}(x)f_{\tau_2}(x)}^{\frac{p}{2}}\\
&\quad + \left(1+|\log\delta|^{-1}\right)^{mp}
    \max_{\substack{ \theta \in \mathcal{P}\bigl(\delta^{\frac{\beta}{2}}\bigr) }} |f_{\widetilde{\theta}}(x)|^p\\
& \leq C_p |\log\delta|^{p} \bigg(
    \sum_{\theta \in \mathcal{P}\bigl(\delta^{\frac{\beta}{2}}\bigr)} |f_\theta(x)|^p
    + \sum_{i=0}^{m-1} \sum_{\tau \in \mathcal{P}(K^{-i})}
      \sum_{\substack{ \tau_1,\tau_2 \in \mathcal{P}(K^{-i-1}) \\
                      \tau_1,\tau_2 \subset \widetilde{\tau} \\
                      \operatorname{dist}(\tau_1,\tau_2) \geq  K^{-i-1} }}
      \abs{f_{\tau_1}(x)f_{\tau_2}(x)}^{\frac{p}{2}}
\bigg).
\end{align*}
Here, the final inequality uses the facts that $m< |\log\delta|$ and $\left(1+|\log\delta|^{-1}\right)^{|\log\delta|}$ is bounded by $e$.
\end{proof}

\end{document}
